% Source: https://artofproblemsolving.com/community/c3h2347509_impossible_chance
% Stable identity: aops-topic-2347509-q01
\textbf{(a)} Suppose every minute had the same probability $p$. If $p=0$, all the probabilities add up to $0$, not $1$. If $p>0$, then already the first $N$ minutes, for any whole number $N>\dfrac1p$, have total probability $Np>1$, which is impossible. Either way we get a contradiction, so the minutes cannot all be equally likely.

\textbf{(b)} The $p_k$ form a geometric series with first term $q$ and ratio $1-q$, where $0<1-q<1$:
\[\sum_{k=1}^{\infty}q(1-q)^{k-1}=\dfrac{q}{1-(1-q)}=1.\]
The probability that the event has not happened in the first $t$ minutes is
\[\sum_{k=t+1}^{\infty}q(1-q)^{k-1}=q(1-q)^t\cdot\dfrac{1}{1-(1-q)}=(1-q)^t.\]
So, given that it has not happened in the first $t$ minutes, the probability that it happens in minute $t+1$ is
\[\dfrac{p_{t+1}}{(1-q)^t}=\dfrac{q(1-q)^t}{(1-q)^t}=\boxed{q},\]
the same for every $t$: in this model, waiting longer neither raises nor lowers the chance for the next minute.

\textbf{(c)} The question has no answer until a probability model for the waiting time is chosen. By (a) there is no model in which all minutes are equally likely, and different models give different answers. In the model of (b), the answer is the same number $q$ for every $t$.
